\lyseis{Α3.1}{Η έννοια του ρητού}
\sol{1} Ο $5$: φυσικός, ακέραιος, ρητός. Ο $-1{,}7$: μόνο ρητός. Ο $\frac{2}{3}$: μόνο ρητός. Ο $-10$: ακέραιος και ρητός. Το $0$: φυσικός, ακέραιος, ρητός. Ο $\frac{12}{4}=3$: φυσικός, ακέραιος, ρητός.
\sol{2} α) Σωστό: $-\frac{6}{2}=-3$. β) Σωστό. γ) Λάθος: ο $\frac{1}{2}$ είναι ρητός αλλά όχι ακέραιος. δ) Σωστό. ε) Λάθος: το $\frac{15}{7}$ είναι το πηλίκο $15:7$. στ) Σωστό: $0=\frac{0}{1}$.
\sol{3} α) $20:4=5$ και $3\cdot5=15$. β) $35:5=7$ και $7\cdot7=49$. γ) $24:6=4$ (και $18:6=3$). δ) $36:4=9$ άρα το $20$ διαιρέθηκε με το $4$: $20:4=5$.
\sol{4} α) $5\cdot9=45=3\cdot15$: ισοδύναμα. β) $4\cdot9=36=6\cdot6$: ισοδύναμα. γ) $16\cdot35=560=7\cdot80$: ισοδύναμα. δ) $21\cdot11=231$ και $13\cdot19=247$: όχι. ε) $14\cdot25=350$ και $15\cdot24=360$: όχι.
\sol{5} Το $\frac{60}{15}=4$ απλοποιείται. Ανάγωγα είναι τα $\frac{24}{35}$ (αφού $24=2^3\cdot3$ και $35=5\cdot7$) και $\frac{19}{6}$. Τα άλλα: $\frac{12}{15}=\frac{4}{5}$ και $\frac{3}{18}=\frac{1}{6}$.
\sol{6} α) $0{,}375$ \quad β) $0{,}35$ \quad γ) $4{,}8$ \quad δ) $-2{,}75$ \quad ε) $0{,}8333\ldots$ \quad στ) $2{,}08$.
\sol{7} α) $\frac{45}{10}=\frac{9}{2}$ \quad β) $\frac{31}{100}$ \quad γ) $\frac{441}{10}$ \quad δ) $-\frac{25}{10}=-\frac{5}{2}$ \quad ε) $\frac{2}{100}=\frac{1}{50}$ \quad στ) $-\frac{1\,125}{1\,000}=-\frac{9}{8}$.
\sol{8} $\mathrm{Α}=-2{,}5=-\frac{5}{2}$, $\mathrm{Β}=-0{,}5=-\frac{1}{2}$, $\Gamma=0{,}5=\frac{1}{2}$, $\Delta=1{,}5=\frac{3}{2}$, $\mathrm{Ε}=2{,}5=\frac{5}{2}$. Κάθε μονάδα έχει διαιρεθεί στη μέση, γι' αυτό όλα τα σημεία είναι «μισά».
\sol{9} Έφαγε $\frac{18}{24}=\frac{18:6}{24:6}=\frac{3}{4}=0{,}75$ της σοκολάτας.
\sol{10} $\frac{45}{60}=\frac{3}{4}$, \ $\frac{20}{60}=\frac{1}{3}$, \ $\frac{40}{60}=\frac{2}{3}$, \ $\frac{12}{60}=\frac{1}{5}$.
\sol{11} Πατάτες $2{,}75=\frac{275}{100}=\frac{11}{4}$ κιλά. Τυρί $0{,}4=\frac{4}{10}=\frac{2}{5}$ του κιλού. Ντομάτες $1{,}25=\frac{125}{100}=\frac{5}{4}$ κιλά.
\sol{12} Κλάσμα $\frac{35}{100}$, δεκαδικός $0{,}35$ και ανάγωγο $\frac{35:5}{100:5}=\frac{7}{20}$.
\sol{13} Φούρνος: $-\frac{1}{2}=-0{,}5$, αριστερά του σπιτιού, $2$ cm αριστερά του $0$. Σχολείο: $\frac{3}{2}=1{,}5$, δεξιά του σπιτιού, $6$ cm δεξιά του $0$. Περίπτερο: $-0{,}25$, αριστερά, $1$ cm αριστερά του $0$. Ο φούρνος και το περίπτερο είναι στην ίδια πλευρά.
\sol{14} Με δεκαδικούς: $\frac{1}{2}=0{,}5$ ενώ $\frac{4}{5}=0{,}8$, άρα δεν είναι ίσα. Στην πίτσα: το μισό δεν είναι ίσο με τα τέσσερα πέμπτα. Σωστά: πολλαπλασιάζω και τους δύο όρους επί $5$: $\frac{1}{2}=\frac{1\cdot5}{2\cdot5}=\frac{5}{10}$.
\sol{15} Αν οι όροι είναι $3k$ και $7k$, η διαφορά τους είναι $4k=20$, άρα $k=5$ και το κλάσμα είναι το $\frac{15}{35}$. Έλεγχος: $15\cdot7=105=35\cdot3$ και $35-15=20$.
\sol{16} $\frac{1}{3}=0{,}333\ldots$ και $\frac{2}{3}=0{,}666\ldots$ Το ψηφίο επαναλαμβάνεται χωρίς τέλος. Δεν γράφεται ως δεκαδικός που τελειώνει, γιατί η διαίρεση $1:3$ δεν δίνει ποτέ υπόλοιπο $0$.
\sol{17} Για παράδειγμα $-\frac{3}{4}=\frac{-3}{4}=\frac{3}{-4}$ και ισοδύναμο $-\frac{6}{8}$ και δεκαδικός $-0{,}75$. Βρίσκεται αριστερά του $0$, στα τρία τέταρτα της διαδρομής προς το $-1$, ανάμεσα στο $-1$ και στο $-0{,}5$.
\sol{18} Ελεύθερη απάντηση. Ο μαθητής πρέπει να γράψει σωστά τα ισοδύναμα, π.χ. $\frac{1}{2}$ φλιτζάνι $=0{,}5=\frac{2}{4}$ και $\frac{3}{4}$ του κιλού $=0{,}75=\frac{6}{8}$.
\sol{19} Ελεύθερη κατασκευή. Οι αρνητικοί και οι αριθμοί με κλασματικό ή δεκαδικό μέρος (όπως $-3$, $0{,}5$, $\frac{1}{3}$) δεν μπαίνουν στους φυσικούς, και όσοι δεν είναι ακέραιοι (όπως $0{,}5$ και $\frac{1}{3}$) δεν μπαίνουν στους ακέραιους. Όλοι μπαίνουν στους ρητούς.
